% Part: computability
% Chapter: recursive-functions
% Section: examples

\documentclass[../../../include/open-logic-section]{subfiles}

\begin{document}

\olfileid{cmp}{rec}{exa}
\olsection{Tuladha Fungsi Rekursif Primitif}

Kita wis nduweni sawetara tuladha fungsi rekursif primitif:
fungsi panjumlahan lan perkalian~$\Add$ lan $\Mult$. Fungsi
identitas $\fn{id}(x) = x$ iku rekursif primitif, amarga
mung~$\Proj{1}{0}$. Fungsi konstan $\fn{const}_n(x) = n$
uga rekursif primitif, amarga bisa ditegesi saka $\Zero$
lan $\Succ$ kanthi komposisi kang ditrapake bola-bali.
Iki migunani nalika kita arep nganggo konstanta ing definisi
rekursif primitif. Umpamane, yen kita arep netepake fungsi
$f(x) = 2 \cdot x$, kita bisa ngasilake kanthi komposisi
saka $\fn{const}_n(x)$ lan perkalian minangka $f(x) =
\Mult(\fn{const}_2(x), \Proj{1}{0}(x))$. Wiwit saiki kita
bakal nggunakake cara iki.

\begin{prop}
  Fungsi pemangkatan $\fn{exp}(x, y) = x^y$ iku rekursif primitif.
\end{prop}

\begin{proof}
  Kita bisa netepake $\fn{exp}$ kanthi rekursi primitif kaya mangkene:
  \begin{align*}
    \fn{exp}(x, 0) & = 1\\
    \fn{exp}(x, y+1) & = \Mult(x, \fn{exp}(x,y)).
    \intertext{Miturut syarat kang tliti, iki durung dadi definisi
      rekursif saka fungsi rekursif primitif. Nanging kanthi
      wangun kang cocog karo aturan, kita nduweni:}
    \fn{exp}(x, 0) & = f(x)\\
    \fn{exp}(x, y+1) & = g(x, y, \fn{exp}(x,y)).
    \intertext{ing ngendi}
    f(x) & = \Succ(\Zero(x)) = 1\\
    g(x, y, z) & = \Mult(\Proj{3}{0}(x, y, z), \Proj{3}{2}(x, y, z)) = x \cdot z
  \end{align*}
  mula $f$ lan $g$ ditegesi saka fungsi rekursif primitif
  kanthi komposisi.
\end{proof}

\begin{prop}
  Fungsi pendahulu $\fn{pred}(y)$ kang ditegesi kanthi
  \[
  \fn{pred}(y) = \begin{cases}
    0 & \text{yen $y=0$}\\
    y-1 & \text{yen ora}
  \end{cases}
  \]
  iku rekursif primitif.
\end{prop}

\begin{proof}
 Gatekna yen
 \begin{align*}
   \fn{pred}(0) & = 0 \text{ lan}\\
   \fn{pred}(y+1) & = y.
 \end{align*}
 Iki meh dadi definisi rekursif primitif. Nanging, miturut
 syarat kang tliti, iki durung cocog karo pola definisi
 kanthi rekursi primitif, amarga pola kasebut mbutuhake
 paling ora siji argumen tambahan~$x$. Uga katon aneh
 amarga $\fn{pred}(y)$ saktemene ora digunakake ing
 definisi $\fn{pred}(y+1)$. Nanging kita bisa netepake dhisik
 $\fn{pred}'(x, y)$ kanthi
 \begin{align*}
   \fn{pred}'(x, 0) & = \Zero(x) = 0,\\
   \fn{pred}'(x, y+1) & = \Proj{3}{1}(x, y, \fn{pred'}(x, y)) = y.
 \end{align*}
lan banjur netepake $\fn{pred}$ saka fungsi kasebut kanthi
komposisi, umpamane minangka
$\fn{pred}(x) = \fn{pred}'(\Zero(x), \Proj{1}{0}(x))$.
\end{proof}

\begin{prop}
  Fungsi faktorial $\fn{fac}(x) = \fact{x} = 1 \cdot 2 \cdot 3
  \cdot \dots \cdot x$ iku rekursif primitif.
\end{prop}

\begin{proof}
  Definisi rekursif primitif kang paling cetha yaiku
  \begin{align*}
    \fn{fac}(0) &= 1\\
    \fn{fac}(y+1) & = \fn{fac}(y) \cdot (y+1).
    \intertext{Miturut aturan, kita kudu luwih dhisik
      netepake fungsi rong panggonan $h$}
    h(x, 0) & = \fn{const}_1(x)\\
    h(x, y+1) & = g(x, y, h(x, y))
    \intertext{ing ngendi $g(x, y, z) = \Mult(\Proj{3}{2}(x, y, z),
      \Succ(\Proj{3}{1}(x, y, z)))$ banjur netepake}
    \fn{fac}(y) & = h(\Proj{1}{0}(y), \Proj{1}{0}(y)) = h(y,y).
  \end{align*}
  Wiwit saiki kita bakal luwih luwes lan ora mesthi menehi
  definisi kang manut aturan kanthi komposisi lan rekursi
  primitif.
\end{proof}

\begin{prop}
  Pangurangan terpotong, $x \tsub y$, kang ditegesi kanthi
  \[
  x \tsub y = \begin{cases}
    0 & \text{yen $x < y$}\\
    x-y & \text{yen ora}
  \end{cases}
  \]
  iku rekursif primitif.
\end{prop}

\begin{proof}
  Kita nduweni:
  \begin{align*}
    x \tsub 0 & = x\\
    x \tsub (y+1) & = \fn{pred}(x \tsub y)
  \end{align*}
\end{proof}

\begin{prop}
  Jarak antarane $x$ lan $y$, $\left|x-y\right|$, iku rekursif primitif.
\end{prop}

\begin{proof}
  Kita nduweni $\left| x-y \right| = (x \tsub y) + (y \tsub x)$,
  mula jarak bisa ditegesi kanthi komposisi saka $+$ lan $\tsub$,
  kang kalorone rekursif primitif.
\end{proof}

\begin{prop}
  Biji maksimum saka $x$ lan $y$, $\fn{max}(x,y)$, iku rekursif primitif.
\end{prop}

\begin{proof}
  Kita bisa netepake $\fn{max}(x,y)$ kanthi komposisi saka $+$ lan $\tsub$:
  \[
  \fn{max}(x,y) \defis x + (y \tsub x).
  \]
  Yen $x$ iku maksimum, yaiku $x \ge y$, mula $y \tsub x = 0$,
  dadi $x + (y \tsub x) = x + 0 = x$. Yen $y$ iku maksimum,
  mula $y \tsub x = y - x$, lan dadi
  $x + (y \tsub x) = x + (y - x) = y$.
\end{proof}

\begin{prop}
  \ollabel{prop:min-pr}
  Biji minimum saka $x$ lan $y$, $\fn{min}(x,y)$, iku rekursif primitif.
\end{prop}

\begin{proof}
  Latihan.
\end{proof}

\begin{prob}
  Buktikna \olref[cmp][rec][exa]{prop:min-pr}.
\end{prob}

\begin{prob}
Tuduhna yen \[
f(x, y) =
2^{(2^{\iddots^{2^{x}}})}\raisebox{1ex}{\bigg\rbrace}
\raisebox{1ex}{\text {$y$ angka $2$}}\] iku rekursif primitif.
\end{prob}

\begin{prob}
Tuduhna yen pambagian wilangan wutuh $d(x, y) = \lfloor x/y \rfloor$
(yaiku pambagian kang ora nggatekake bagean sawise tandha desimal)
iku rekursif primitif. Yen $y = 0$, kita netepake
$d(x, y) = 0$. Wenehana definisi eksplisit kanggo~$d$
nganggo rekursi primitif lan komposisi.
\end{prob}

\begin{prop}
Himpunan fungsi rekursif primitif katutup tumrap
rong operasi iki:
\begin{enumerate}
\item Jumlah winates: yen $f(\vec x, z)$ iku rekursif primitif,
mula fungsi
\[
g(\vec x, y) \defis \sum_{z = 0}^y f(\vec x, z).
\]
uga rekursif primitif.
\item Asil kali winates: yen $f(\vec x, z)$ iku rekursif
primitif, mula fungsi
\[
h(\vec x, y) \defis \prod_{z = 0}^y f(\vec x, z).
\]
uga rekursif primitif.
\end{enumerate}
\end{prop}

\begin{proof}
Umpamane, jumlah winates ditegesi kanthi rekursif liwat
persamaan iki:
\begin{align*}
  g(\vec x, 0) & = f(\vec x, 0)\\
  g(\vec x, y+1) & = g(\vec x, y) + f(\vec x, y+1).
\end{align*}
\end{proof}

\end{document}
